% Paper 2, Section 6 — The datapath-degradation principle (evidence E4).
% Numbers verified against ledger entries bb79f7f (char-LM control), 5bb7bcc (arm
% mis-assignment discovered; quantizer contrast identified), 8d050cb (M=16 n=3), 48510c8
% (M=32 n=3, pre-registered floor met).
% CLAIMS DISCIPLINE: the quantizer contrast is POST-HOC-IDENTIFIED and must carry that label;
% the pre-registered arm as written tested baseline tuning, not datapaths. The noise-benefit
% trend across loads runs TOWARD char-LM and must be reported as weakening the strength line.

\section{The datapath-degradation principle}
\label{sec:principle}

\subsection{Statement}
\label{sec:principle-statement}

Sections~\ref{sec:charlm} and~\ref{sec:copy} give matched-communication-rate variant
orderings on two workloads, and they invert. One principle predicts both:

\begin{quote}
\emph{A neuromorphic SSM datapath is cheap exactly when it degrades the part of the
computation the workload does not depend on.} Statistical sequence workloads need an exact
graded \emph{output} distribution and tolerate a lossy state, so the analog-state route is
the best neuromorphic option there. Precise-recall workloads need an exact recurrent
\emph{state} and tolerate a spiked output, so the output-spiking route is the only
neuromorphic option that survives there.
\end{quote}

\subsection{Why state fidelity alone is not the axis}
\label{sec:principle-fidelity}

A simpler candidate explanation --- rank the variants by how exact their recurrent state is
--- fits the copy ordering perfectly: at matched emission, margin kept falls monotonically
with state fidelity, continuous ($0.87$) $>$ lossy-analog ($0.50$) $>$ one-bit ($0.074$). But
it fails on char-LM: the output-spiking variant has the \emph{most} exact recurrent state on
both tasks (state activity exactly $1.0000 \pm 0.0000$ in every cell of both tables, the most
reproducible measurement in this study) yet is the \emph{worst} variant on char-LM and the
best on copy. State fidelity predicts one table; only the two-sided principle --- which part
of the datapath does the task read? --- predicts both.

\subsection{An out-of-sample test: the quantizer contrast}
\label{sec:principle-quantizer}

The principle makes a falsifiable prediction about any single degradation applied to the
\emph{same} datapath on both workloads: an inference-time loss of state precision should be
near-free on the statistical task and costly on the precise-recall task. The regularization
control of Section~\ref{sec:controls} contains exactly such a degradation: the 6-bit state
quantizer (unlike the additive noise, which is training-time only, the quantizer acts at
inference). Applied to the \emph{digital} baseline at matched noise $\sigma{=}0.02$, the same
quantizer costs:

\begin{center}
\footnotesize
\setlength{\tabcolsep}{4pt}
\begin{tabular}{l r l}
\toprule
Workload & $\Delta$bpc & state footprint \\
\midrule
char-LM         & $-0.002$ & --- \\
copy $M{=}16$   & $+0.369$ & $1.000{\to}0.971{\pm}0.015$ \\
copy $M{=}32$   & $+0.385$ & $1.000{\to}0.976{\pm}0.005$ \\
\bottomrule
\end{tabular}
\end{center}

\noindent
Both copy rows are three seeds, and $\Delta$bpc is the quantizer-alone cost at matched
$\sigma{=}0.02$. The identical hardware degradation is free on the statistical workload and costs
$\sim 0.37$--$0.39$ bpc on the precise-recall workload --- the principle's prediction, out of
sample, on a datapath (plain digital) that none of the four variants' comparisons touched.
The effect is load-robust: a floor of ``$\ge +0.37$ at $M{=}32$, smaller counts as evidence
against'' was pre-registered after the $M{=}16$ result and met ($+0.385$ vs.\ $+0.369$; flat
to slightly rising, within arm sd $\sim 0.04$--$0.05$ --- we say ``does not attenuate with
load,'' not ``grows''). Quantization alone still keeps $0.78$ of the digital margin at
$M{=}32$, far above spikeout's $0.42$ and analog's $0.28$, so state-precision loss explains
part but not most of the neuromorphic copy deficit; the event and spike mechanisms add their
own cost.

\paragraph{Honesty label: this test is post-hoc-identified.}
The control row was pre-registered with the \emph{noise} arm bound to the principle's
prediction, and that assignment was wrong: inspection of the implementation showed the
additive noise is gated behind training (\texttt{if self.training}) while the quantizer is
not, so the noise arm is a training-time regularizer that tests baseline tuning, not an
inference datapath degradation --- and its landing (noise \emph{helps} on copy too) says
nothing about the principle. The quantizer arm is the row's only inference-time state
degradation, and it was identified as the real test only after the arms landed. We therefore
report the contrast as post-hoc but mechanistically motivated, with one genuinely
pre-registered element (the $M{=}32$ floor, set before those cells finished). The
pre-registered test \emph{as written} failed by mis-specification.

\subsection{What weakens the principle}
\label{sec:principle-against}

Two observations run against a strong reading. First, the training-time noise benefit does
not separate the workloads the way a naive reading would hope: it is $-0.150 \pm 0.033$ bpc
at copy $M{=}16$, $-0.215 \pm 0.011$ at $M{=}32$, and $-0.309 \pm 0.024$ on char-LM --- the
trend across loads runs \emph{toward} the char-LM value, so ``the regularizer helps the
statistical task more'' is at best a weak supporting line and we do not lean on it. Second,
the principle has been tested out of sample exactly once (the quantizer contrast), on two
tasks and one architecture family; Section~\ref{sec:limitations} states the scale caveats.

Third --- and decisive for how the principle may be used --- when employed as a
\emph{forecast} on settings it had not seen, the principle is zero for two. It predicted
that an event-driven readout, being an output degradation, would be cheap on the
precise-recall task, where output-spiking was already the best variant; the gate did not
degrade recall, it eliminated it (Section~\ref{sec:hardware-stranded}). It then predicted
that a lossy analog state, being a state degradation, would be cheap on a streaming
statistical perception task; the measured deficit is $10$--$19$ accuracy points at every
threshold, robust to a $3\times$ budget increase (Section~\ref{sec:hardware-shape}). The
principle correctly organizes every degradation this paper measured, and both failed
forecasts have identifiable causes in hindsight (staleness is not the LIF degradation
axis; ``statistical'' is too coarse a workload class) --- but hindsight is the point. We
therefore state it throughout as a post-hoc ordering of measured degradations, not a
predictive law, and the designer's rule it supports is to \emph{measure} which side of
the datapath the workload's loss depends on, not to predict it.
